\documentclass{oicontest-v2}

\title{NOI2024 题面演示}
\subtitle{Day 1}
\author{ExplodingKonjac}
\date{\today}

\listadd{\oicattention}{选手可使用快捷启动页面中的工具 selfEval 进行自测。在将待测程序（不必是全部题目）放到题目目录下后，即可选择全部或部分题目进行自测。注意：自测有次数限制，且自测结果仅用于选手调试，并不做为最终正式成绩。}

\addproblem{集合}{set}{传统型}{2.0}{512}{20}{是}
\addproblem{百万富翁}{richest}{交互型}{6.0}{512}{2}{否}
\addproblem{树的定向}{tree}{传统型}{3.0}{2048}{25}{是}

\begin{document}

\oictitlepage

\oicnextproblem % set

\oicstatement

小 Y 和小 S 在玩一个游戏。

给定正整数 $m$，定义\textbf{基本集合}为大小为 $3$，元素在 $1\sim m$ 内的集合。例如：给定 $m=4$，则集合 $\{1,2,3\}$ 与集合 $\{2,3,4\}$ 都是基本集合。

定义\textbf{集合序列}为由基本集合构成的序列，例如，$A=[\{1,2,3\},\{2,3,4\}]$ 是一个集合序列，其中 $A[1]=\{1,2,3\}$，$A[2]=\{2,3,4\}$ 都是基本集合。

对于一个 $1\sim m$ 的排列 $p[1],p[2],\dots,p[m]$ 与集合 $S\subseteq \{1,2,\dots,m\}$，定义 $f_p(S)$ 为将 $S$ 内每一个元素 $x$ 置换为 $p[x]$ 后所得到的集合，即 $f_p(S)=\{p[x]|x\in S\}$。

对于两个长度为 $k$ 的集合序列 $A,B$，定义 $A$ 和 $B$ \textbf{等价}当且仅当存在一个 $1\sim m$ 的排列 $p$，使得 $A$ 置换排列 $p$ 后得到 $B$，即对于所有 $1\leq i\leq k$，$f_p(A[i])=B[i]$。

给定两个长度为 $n$ 的集合序列 $A,B$，有 $q$ 次询问。每次小 S 会询问小 Y，在给定 $l,r$ 的情况下，判断集合序列 $[A[l],A[l+1],\dots,A[r]]$ 与集合序列 $[B[l],B[l+1],\dots,B[r]]$ 是否等价。

时光荏苒，小 S 和小 Y 也会散去。而我们和一个人保持连接的方式就是记住，仅此而已。

\oicinputformat

输入的第一行包含三个正整数 $n,m,q$，分别表示集合序列的长度，元素范围和询问次数。

输入的第二行包含 $3n$ 个正整数，第 $3i-2,3i-1,3i\ (1\leq i\leq n)$ 个正整数分别表示 $A[i]$ 的三个元素。保证这三个元素均在 $[1,m]$ 范围内且互不相同。

输入的第三行包含 $3n$ 个正整数，第 $3i-2,3i-1,3i\ (1\leq i\leq n)$ 个正整数分别表示 $B[i]$ 的三个元素。保证这三个元素均在 $[1,m]$ 范围内且互不相同。

接下来 $q$ 行，每行包含两个正整数 $l,r$，表示一次询问。

\oicoutputformat

输出 $q$ 行，每行包含一个字符串 $\tt{Yes}$ 或 $\tt{No}$，表示对应询问的两个序列是否等价。

\oicsampleinput{1}
\begin{oicsample}
4 4 10
1 2 3 1 2 3 1 2 4 1 2 3
1 2 4 2 3 4 1 2 3 2 3 4
1 1
1 2
1 3
1 4
2 2
2 3
2 4
3 3
3 4
4 4
\end{oicsample}

\oicsampleoutput{1}
\begin{oicsample}
Yes
No
No
No
Yes
Yes
Yes
Yes
Yes
Yes
\end{oicsample}

\oicsampleexplain{1}

以下用 $(l,r)$ 表示对 $l,r$ 的询问：

\begin{enumerate}
	\item 对于询问 $(1,1)$，令排列 $p=[1,2,4,3]$，则 $f_p(A_1)=\{p[1],p[2],p[3]\}=\{1,2,4\}=B[1]$，因此该询问对应的两个序列等价。
	\item 对于询问 $(1,2),(1,3),(1,4)$，由于 $A[1]=A[2]$ 但 $B[1]\neq B[2]$，因此这些询问对应的两个序列都不等价。
	\item 对于询问 $(2,2),(2,3),(2,4),(3,3),(3,4),(4,4)$，令排列 $p=[2,3,4,1]$，则 $f_p(A_2)=\{p[1],p[2],p[3]\}=\{2,3,4\}=B_2$，$f_p(A_3)=\{p[1],p[2],p[4]\}=\{1,2,3\}=B_3$，$f_p(A_4)=\{p[1],p[2],p[3]\}=\{2,3,4\}=B_4$，因此这些询问对应的两个序列都等价。 
\end{enumerate}

\oicsamplefile{2}

这个样例满足测试点 $1\sim 3$ 的约束条件。

\oicsamplefile{3}

这个样例满足测试点 8 的约束条件。

\oicsamplefile{4}

这个样例满足测试点 15、16 的约束条件。

\oicconstraints

对于所有测试数据保证：$1\leq n\leq 2\times 10^5$，$3\leq m\leq 6\times 10^5$，$1\leq q\leq 10^6$，$1\leq l\leq r\leq n$

\begin{oictable}{cccc}
	测试点编号 & $n\leq$ & $m\leq$ & $q\leq$ \\
	$1\sim 3$ & $50$ & $4$ & $50$ \\
	$4\sim 6$ & $50$ & $5$ & $50$ \\
	$7$ & $200$ & $4$ & $200$ \\
	$8$ & $200$ & $5$ & $200$ \\
	$9$ & $200$ & $4$ & $2\times 10^5$ \\
	$10$ & $200$ & $5$ & $2\times 10^5$ \\
	$11$ & $2\times 10^5$ & $4$ & $2\times 10^5$ \\
	$12$ & $2\times 10^5$ & $5$ & $2\times 10^5$ \\
	$13,14$ & $2\,000$ & $6\,000$ & $10^3$ \\
	$15,16$ & $2\,000$ & $6\,000$ & $10^6$ \\
	$17,18$ & $2\times 10^4$ & $6\times 10^4$ & $10^2$ \\
	$19,20$ & $2\times 10^5$ & $6\times 10^5$ & $10^6$ \\
\end{oictable}

\oicnextproblem % richest

\textbf{这是一道交互题。}\vspace{-1em}

\oicstatement

小 Y 的银行有 $N$ 个客户，编号为 $0$ 到 $N-1$。客户 $i$ 有 $W_i$ 元存款，且\textbf{客户之间的存款金额互不相同}。

小 P 是小 Y 的深度合作伙伴，他希望知道哪个客户的存款最多。小 P 无法直接获取客户的存款金额，但他可以依次发送若干次\textbf{请求}，每次请求包含若干个\textbf{查询}，每个查询是一个二元组 $(i,j)$，表示小 P 想知道客户 $i$ 和客户 $j$ 的存款金额哪个更多。如果 $W_i>W_j$，小 Y 会回答 $i$，否则回答 $j$。

小 P 的\textbf{请求数} $t$ 和所有请求的\textbf{查询}次数总和 $s$ 有上限，他希望你帮他写一个程序，来找到存款最多的客户。

\oicimplementation

\textbf{选手不需要，也不应该实现 \texttt{main} 函数。}

选手应确保提交的程序包含头文件 \texttt{richest.h}，可在程序开头加入以下代码实现：

\begin{minted}{cpp}
#include "richest.h"
\end{minted}

选手需要实现以下函数：

\begin{minted}{cpp}
int richest(int N,int T,int S);
\end{minted}

\begin{itemize}
	\item  $N$ 表示客户的数量；
	\item $T$ 表示对于当前函数调用，\textbf{请求}数 $t$ 不应超过此值；
	\item $S$ 表示对于当前函数调用，所有请求的\textbf{查询}次数总和 $s$ 不应超过此值；
	\item 该函数需要返回存款最多的客户的编号；
	\item 对于每个测试点，该函数\textbf{会被交互库调用恰好 $10$ 次}。
\end{itemize}

选手可以通过调用以下函数向交互库发送一次\textbf{请求}：

\begin{minted}{cpp}
std::vector<int> ask(std::vector<int> a, std::vector<int> b);
\end{minted}

\begin{itemize}
	\item 在调用 \texttt{ask} 函数时需要保证传入参数 $a$ 和 $b$ 的长度相同，且其中的每个元素都必须是小于 $N$ 的非负整数，表示该\textbf{请求}中的所有\textbf{查询}。
	\item 该函数会返回一个类型为 \texttt{std::vector <int>} 且长度与 $a$ 和 $b$ 相同的变量，设为 $c$，其中 $c[i]$ 表示在客户 $a[i]$ 和 $b[i]$ 中存款较多的客户的编号。
\end{itemize}

题目保证在规定的\textbf{请求}和\textbf{查询}次数限制下，交互库运行的时间不超过 $3$ 秒，交互库使用的内存大小固定，且不超过 $256$ MiB。

\oictesting

\textbf{试题目录下的 \texttt{grader.cpp} 是提供的交互库参考实现，最终测试时所用的交互库实现与该参考实现有所不同，因此选手的解法不应该依赖交互库实现。}

选手可以在本题目录下使用如下命令编译得到可执行程序：

\begin{minted}{text}
g++ grader.cpp richest.cpp -o richest -O2 -std=c++14 -static
\end{minted}

对于编译得到的可执行程序：

\begin{itemize}
	\item 可执行文件将从\textbf{标准输入}读入以下格式的数据：
	\begin{itemize}
		\item 输入的第一行包含四个非负整数 $N,T,S,R$，其中 $R$ 是交互库生成测试数据的随机种子。
	\end{itemize}
	\item  输入完成后，交互库将调用 $10$ 次函数 \texttt{richest}，用输入的参数生成的测试数据进行测试。\texttt{richest} 函数返回后，交互库会输出以下信息：
	\begin{itemize}
		\item 输出的前 $10$ 行中，每行首先包含三个整数 $r,t,s$，表示该次执行的结果。其中 $r$ 是 \texttt{richest} 函数的返回值，$t$ 和 $s$ 的含义如题目描述中所示，然后包含该次运行的正确性等信息。
		\item 输出的第 $11$ 行包含 $10$ 次运行的总信息。
	\end{itemize}
\end{itemize}

\subsection{交互示例}

假设可执行文件生成的测试数据为 $N=4$，$W=[101,103,102,100]$，$T=100$，$S=100$。

下面是一个正确的交互过程：

\begin{oictable}{ccc}
	选手程序 & 交互库 & 说明 \\
	~ & 调用 richest(4, 100, 100) & 开始测试 \\
	调用 ask([0, 2], [1, 3]) & 返回 $[1,2]$ & $W_0<W_1$，$W_2>W_3$ \\
	调用 ask([0, 2, 3], [1, 1, 1]) & 返回 $[1,1,1]$ & $W_0<W_1$，$W_2<W_1$，$W_3<W_1$ \\
	运行结束并返回 $1$ & 向屏幕打印交互结果 & 交互结束，结果正确 \\
\end{oictable}

在这个例子中，$r=1$，$t=2$，$s=5$，满足请求与查询次数的限制。

\subsection{下发文件说明}

在本试题目录下：

\begin{enumerate}
	\item \texttt{grader.cpp} 是提供的交互库参考实现。
	\item \texttt{richest.h} 是头文件，选手不用关心具体内容。
	\item \texttt{template\_richest.cpp} 是提供的示例代码，选手可在此代码的基础上实现。
\end{enumerate}

\textbf{选手注意对所有下发文件做好备份，最终评测时只测试本试题目录下的 \texttt{richest.cpp}，对该程序以外文件的修改不会影响评测结果。}

\oicconstraints

对于所有测试数据保证：所有 $W_i$ 两两不同。

本题共 $2$ 个测试点，每个测试点的分值和数据范围见下表。

\begin{oictable}{ccccc}
	测试点编号 & 分值 & $N=$ & $T=$ & $S=$ \\
	$1$ & $15$ & $1\,000$ & $1$ & $499\,500$ \\
	$2$ & $85$ & $1\,000\,000$ & $20$ & $2\,000\,000$ \\
\end{oictable}

\oicscoring

注意：

\begin{itemize}
	\item 选手不应通过非法方式获取交互库的内部信息，如试图直接读取数组 $W$ 的值，或直接与标准输入、输出流进行交互。此类行为将被视为作弊。
	\item \textbf{最终的评测交互库与样例交互库的实现不同，且可能是适应性的：在不与 \texttt{ask} 此前返回的结果相矛盾的前提下，最终的评测交互库可能会动态调整 $W$ 的值。}
\end{itemize}

\textbf{本题首先会受到和传统相同的限制}。例如编译错误会导致整道题目得 $0$ 分，运行时错误、超过时间限制、超过空间限制都会导致相应测试点得 0 分。选手只能在程序中访问自己定义的和交互库给出的变量或数据，及其相应的内存空间。尝试访问其他位置空间将可能导致编译错误或运行错误。

在每次 \texttt{richest} 函数调用中，程序使用的请求次数 $t$ 和所有请求的查询次数总和 $s$ 需在对应限制下，否则将会获得 $0$ 分。

在上述条件基础上：

\begin{itemize}
	\item 在测试点 $1$ 中，程序得到满分当且仅当 \texttt{ask} 函数调用合法且 \texttt{richest} 函数返回的答案正确；
	\item 在测试点 $2$ 中，程序得到的分数将按照以下方式计算：
	\begin{itemize}
		\item 若 \texttt{ask} 函数调用不合法，则获得 $0$ 分；
		\item 若 \texttt{ask} 函数调用均合法，设 $\max t$ 表示多次调用 \texttt{richest} 函数所得的 $t$ 的最大值，$\max s$ 表示 $s$ 的最大值，则程序将获得 $\lfloor 85 \cdot f(\max t) \cdot g(\max s)\rfloor$ 分，其中 $f$ 与 $g$ 的计算方式如下表所示：
	\end{itemize}
\end{itemize}

\begin{oictable}{cc}
	$\max t$ & $f(\max t)$ \\
	$\max t\leq 8$ & $1$ \\
	$9\leq \max t\leq 20$ & $1-\frac{1}{4}\sqrt{\max t-8}$ \\
\end{oictable}

\begin{oictable}{cc}
	$\max s$ & $g(\max s)$ \\
	$\max s\leq 1\,099\,944$ & $1$ \\
	$1\,099\,945\leq \max s\leq 1\,100\,043$ & $1-\frac{1}{6} \log_{10} (\max s-1\,099\,943)$ \\
	$1\,100\,044\leq \max s\leq 2\,000\,000$ & $\frac{2}{3}-\frac{1}{1\,500}\sqrt{\max s-1\,100\,043}$ \\
\end{oictable}

以下是测试点 $2$ 中，不同的 $t$ 和 $s$ 对得分影响的示例。

\begin{oictable}{ccc}
	$\max t$ & $\max s$ & 测试点 $2$ 的得分 \\
	$=20$ & \SetCell[r=12]{c}$\le 1\,099\,944$ & $11$ \\
	$=19$ & & $14$ \\
	$=18$ & & $17$ \\
	$=17$ & & $21$ \\
	$=16$ & & $24$ \\
	$=15$ & & $28$ \\
	$=14$ & & $32$ \\
	$=13$ & & $37$ \\
	$=12$ & & $42$ \\
	$=11$ & & $48$ \\
	$=10$ & & $54$ \\
	$=9$ & & $63$ \\
	\SetCell[r=16]{c}$\le 8$ & $\in [1\,099\,974,1\,099\,978]$ & $63$ \\
	& $\in [1\,099\,969,1\,099\,973]$ & $64$ \\
	& $\in [1\,099\,965,1\,099\,968]$ & $65$ \\
	& $\in [1\,099\,962,1\,099\,964]$ & $66$ \\
	& $\in [1\,099\,959,1\,099\,961]$ & $67$ \\
	& $\in [1\,099\,957,1\,099\,958]$ & $68$ \\
	& $\in [1\,099\,955,1\,099\,956]$ & $69$ \\
	& $\in [1\,099\,953,1\,099\,954]$ & $70$ \\
	& $=1\,099\,952$ & $71$ \\
	& $=1\,099\,951$ & $72$ \\
	& $\in [1\,099\,949,1\,099\,950]$ & $73$ \\
	& $=1\,099\,948$ & $75$ \\
	& $=1\,099\,947$ & $76$ \\
	& $=1\,099\,946$ & $78$ \\
	& $=1\,099\,945$ & $80$ \\
	& $\le 1\,099\,944$ & $85$ \\
\end{oictable}

\oicnextproblem % tree

\oicstatement

给定一棵含有 $n$ 个顶点的树，顶点从 $1$ 到 $n$ 编号，树上第 $i$（$1\leq i\leq n-1$） 条边连接顶点 $u_i$ 和 $v_i$。

现在，我们想要给树的每条边一个定向。任何一个定向都可以用一个长度为 $n-1$ 的字符串 $S=s_1s_2\ldots s_{n-1}$ 来描述。其中 $s_i=0$ 代表第 $i$ 条边定向为 $u_i \to v_i$，否则 $s_i=1$ 代表第 $i$ 条边定向为 $v_i\to u_i$。

给定 $m$ 个顶点对 $(a_i,b_i)$，其中 $1\leq a_i,b_i \leq n$ 且 $a_i\neq b_i$。

一个\textbf{完美定向}定义为：在此定向下，对于任意 $1\leq i\leq m$，$a_i$ \textbf{不能到达} $b_i$。

试求在所有完美定向中，所对应的字符串字典序最小的定向。\textbf{数据保证存在至少一个完美定向}。

定义字符串 $S=s_1\ldots s_{n-1}$ 的字典序小于 $T=t_1\ldots t_{n-1}$：若存在一个下标 $k$，使得 $s_1=t_1,\ldots,s_{k-1}=t_{k-1},s_k<t_k$。

\oicinputformat

输入的第一行包含三个非负整数 $c,n,m$，分别表示测试点编号，树的点数，顶点对的个数。$c=0$ 表示该测试点为样例。

接下来 $n-1$ 行，每行包含两个正整数 $u_i,v_i$ 表示树的一条边。保证 $1\leq u_i,v_i\leq n$ 且这 $n-1$ 条边构成了一棵树。

接下来 $m$ 行，每行包含两个正整数 $a_i,b_i$。保证 $1\leq a_i,b_i \leq n$ 且 $a_i\neq b_i$。

\oicoutputformat

输出一行包含一个字符串 $S=s_1\ldots s_{n-1}$，表示字典序最小的完美定向所对应的 $01$ 字符串。

\oicsampleinput{1}
\begin{oicsample}
0 4 2
1 2
2 3
3 4
3 2
1 4
\end{oicsample}

\oicsampleoutput{1}
\begin{oicsample}
001
\end{oicsample}

\oicsampleexplain{1}

在该样例中，若 $S=000$，则该定向中 $1$ 能到达 $4$（存在路径 $1\to 2\to 3\to 4$），因而不是完美定向。若 $S=001$，则该定向中 $3$ 不能到达 $2$，$1$ 不能到达 $4$，因面是完美定向。故答案为 $001$。

\oicsampleinput{2}
\begin{oicsample}
0 6 8
5 1
2 3
1 2
5 6
4 3
4 3
5 1
6 3
5 4
1 4
5 2
3 6
6 2
\end{oicsample}

\oicsampleoutput{2}
\begin{oicsample}
10101
\end{oicsample}

\oicsampleexplain{2}

在该样例中，一组完美定向必定满足 $4$ 不能到达 $3$，$5$ 不能到达 $1$。故 $s_1=s_5=1$。若 $s_2=s_3=0$，则存在路径 $1\to 2\to 3\to 4$，故 $1$ 可到达 $4$。故其不是完美定向。因此，所有完美定向必定满足 $S$ 的字典序不小于 $10101$。且容易验证 $S=10101$ 时，对应的定向是完美定向。

\oicsamplefile{3}

这个样例满足测试点 $1\sim 3$ 的约束条件。

\oicsamplefile{4}

这个样例满足测试点 $4\sim 6$ 的约束条件。

\oicsamplefile{5}

这个样例满足测试点 $7,8$ 的约束条件。

\oicsamplefile{6}

这个样例满足测试点 $9,10$ 的约束条件。

\oicconstraints

对于所有测试数据保证 $2\leq n\leq 5\times 10^5$，$1\leq m\leq 5\times 10^5$，$1\leq u_i,v_i\leq n$ 且所有的边构成了一棵树，$1\leq a_i,b_i \leq n$ 且 $a_i\neq b_i$。

数据保证存在至少一个完美定向。

\begin{oictable}{cccc}
	测试点编号 & $n$ & $m$ & 特殊性质 \\
	$1\sim 3$ & $\leq 15$ & $\leq 50$ & 无 \\
	$4\sim 6$ & $\leq 300$ & $\leq 300$ & 无 \\
	$7,8$ & $\leq 400$ & $=(n-1)(n-2)$ & A \\
	$9,10$ & $\leq 2\,000$ & $\leq 2\,000$ & B \\
	$11\sim 14$ & $\leq 2\,000$ & $\leq 2\,000$ & 无 \\
	$15,16$ & $\leq 10^5$ & $\leq 10^5$ & B \\
	$17,18$ & $\leq 10^5$ & $\leq 10^5$ & 无 \\
	$19\sim 21$ & $\leq 2\times 10^5$ & $\leq 2\times 10^5$ & 无 \\
	$22\sim 25$ & $\leq 5\times 10^5$ & $\leq 5\times 10^5$ & 无 \\
\end{oictable}

特殊性质 A：保证 $(a,b)$ 出现在 $(a_i,b_i)$ 中当且仅当 $a\neq b$ 且 $a,b$ 在树上不相邻。

特殊性质 B：保证树上编号为 $1$ 的顶点与其他每个顶点均相邻。

\end{document}